\section*{NeurIPS Paper Checklist}

\begin{enumerate}

\Needspace{8\baselineskip}

\item {\bf Claims}
    \item[] Question: Do the main claims made in the abstract and introduction accurately reflect the paper's contributions and scope?
    \item[] Answer: \answerYes{}
    \item[] Justification: Sections~\ref{sec:introduction}--\ref{sec:experiments} distinguish conditional classification bounds, operational information targets, and empirical sensitivity findings. Section~\ref{sec:limitations} states the limits on generalization, nuisance attribution, and behavioral repair.

\Needspace{8\baselineskip}

\item {\bf Limitations}
    \item[] Question: Does the paper discuss the limitations of the work performed by the authors?
    \item[] Answer: \answerYes{}
    \item[] Justification: Section~\ref{sec:limitations} discusses specification dependence, local-theory assumptions, nominal audit coverage, development history, benchmark scope, and the distinction between information gains and predictive improvement.

\Needspace{8\baselineskip}

\item {\bf Theory assumptions and proofs}
    \item[] Question: For each theoretical result, does the paper provide the full set of assumptions and a complete (and correct) proof?
    \item[] Answer: \answerYes{}
    \item[] Justification: The theoretical statements give their assumptions, with proofs in Appendices~\ref{app:core-proofs}, \ref{app:estimation-certification}, \ref{app:hidden-slice-discovery}, and~\ref{app:information-interventions}. The classification result explicitly retains distinguishability, decision-separation, and local-radius conditions.

    \Needspace{8\baselineskip}

    \item {\bf Experimental result reproducibility}
    \item[] Question: Does the paper fully disclose all the information needed to reproduce the main experimental results of the paper to the extent that it affects the main claims and/or conclusions of the paper (regardless of whether the code and data are provided or not)?
    \item[] Answer: \answerYes{}
    \item[] Justification: Appendices~\ref{app:estimation-certification}--\ref{app:information-interventions} specify the data splits, models, objectives, hyperparameters, seed sets, selection rules, interventions, and evaluation procedures. The public release provides training and evaluation code, experiment-specific rerun instructions, and numerical replay of the retained main-result artifacts. Independent reruns are distinguished from exact replay of historical outputs that were not retained.


\Needspace{8\baselineskip}


\item {\bf Open access to data and code}
    \item[] Question: Does the paper provide open access to the data and code, with sufficient instructions to faithfully reproduce the main experimental results, as described in supplemental material?
    \item[] Answer: \answerYes{}
    \item[] Justification: The public \href{https://github.com/E-zClap/fisher-glass-learning/tree/camera-ready-v14}{\texttt{camera-ready-v14} release} provides the experiment code, dependencies, data-access and preparation instructions, commands for training and evaluation, and retained numerical inputs and outputs. Its reproduction guides map the main results to these workflows and identify historical outputs unavailable for exact replay. Public benchmark data and pretrained weights are obtained from their original providers.


\Needspace{8\baselineskip}


\item {\bf Experimental setting/details}
    \item[] Question: Does the paper specify all the training and test details (e.g., data splits, hyperparameters, how they were chosen, type of optimizer) necessary to understand the results?
    \item[] Answer: \answerYes{}
    \item[] Justification: Appendices~\ref{app:hidden-slice-discovery}--\ref{app:information-interventions} give the data roles, fitting and calibration procedures, optimizers, hyperparameters, seed sets, and selection rules. They distinguish original cohorts, fresh audit replications, and post hoc reanalyses.

\Needspace{8\baselineskip}

\item {\bf Experiment statistical significance}
    \item[] Question: Does the paper report error bars suitably and correctly defined or other appropriate information about the statistical significance of the experiments?
    \item[] Answer: \answerYes{}
    \item[] Justification: Appendix~\ref{app:experimental-conventions} defines nominal normal-approximation seed intervals and distinguishes Student-$t$ intervals, reported standard deviations, and conditional bootstrap uncertainty; Appendix~\ref{app:implemented-audit} separately qualifies nominal Wald values. Acquisition tables explicitly report means without intervals because their per-seed outputs are unavailable.

\Needspace{8\baselineskip}

\item {\bf Experiments compute resources}
    \item[] Question: For each experiment, does the paper provide sufficient information on the computer resources (type of compute workers, memory, time of execution) needed to reproduce the experiments?
    \item[] Answer: \answerNo{}
    \item[] Justification: Appendices~\ref{app:experimental-details}, \ref{app:regularized-training}, and~\ref{app:extended-training} identify the reported CPU/GPU hardware and available memory and timing information. Per-run execution times, peak memory, and total compute including development runs were not recorded comprehensively.

\Needspace{8\baselineskip}

\item {\bf Code of ethics}
    \item[] Question: Does the research conducted in the paper conform, in every respect, with the NeurIPS Code of Ethics \url{https://neurips.cc/public/EthicsGuidelines}?
    \item[] Answer: \answerYes{}
    \item[] Justification: The work uses synthetic data and existing research benchmarks, with no participant recruitment. The release includes selected figure illustrations but excludes full benchmark datasets and pretrained weights. Section~\ref{sec:limitations} and the experimental appendices disclose the scope of the evidence and distinguish approximate numerical audits from formal guarantees.


\Needspace{8\baselineskip}


\item {\bf Broader impacts}
    \item[] Question: Does the paper discuss both potential positive societal impacts and negative societal impacts of the work performed?
    \item[] Answer: \answerYes{}
    \item[] Justification: The audits may help identify subpopulations vulnerable to nuisance changes, but treating a restricted audit as a universal guarantee could create false assurance and obscure harms in unaudited populations. Section~\ref{sec:limitations} limits conclusions to declared populations and specifications and calls for prospective validation of deployment objectives.

\Needspace{8\baselineskip}

\item {\bf Safeguards}
    \item[] Question: Does the paper describe safeguards that have been put in place for responsible release of data or models that have a high risk for misuse (e.g., pre-trained language models, image generators, or scraped datasets)?
    \item[] Answer: \answerNA{}
    \item[] Justification: The work does not release models or datasets with a high risk of misuse. The accompanying research code and numerical artifacts exclude full benchmark datasets and pretrained model weights.

\Needspace{8\baselineskip}

\item {\bf Licenses for existing assets}
    \item[] Question: Are the creators or original owners of assets (e.g., code, data, models), used in the paper, properly credited and are the license and terms of use explicitly mentioned and properly respected?
    \item[] Answer: \answerNo{}
    % ASSET_LICENSE_SCOPE: No reflects incomplete historical license/version records; excluding external assets does not establish a new license grant.
    \item[] Justification: The original datasets, models, and methods are credited, and \path{DATA_LICENSES.md} links upstream terms, but complete license and version records for every historical dependency were not retained. The release excludes full benchmark datasets and pretrained weights; selected figure illustrations retain their source attribution and do not receive a new license.

\Needspace{8\baselineskip}

\item {\bf New assets}
    \item[] Question: Are new assets introduced in the paper well documented and is the documentation provided alongside the assets?
    \item[] Answer: \answerYes{}
    \item[] Justification: The accompanying code and numerical-artifact package provides a README, provenance and hash manifests, replay commands, and explicit descriptions of unavailable historical outputs. Appendices~\ref{app:experimental-details} and~\ref{app:information-interventions} identify the implemented protocols and evidence supporting the reported results.

\Needspace{8\baselineskip}

\item {\bf Crowdsourcing and research with human subjects}
    \item[] Question: For crowdsourcing experiments and research with human subjects, does the paper include the full text of instructions given to participants and screenshots, if applicable, as well as details about compensation (if any)?
    \item[] Answer: \answerNA{}
    \item[] Justification: The research does not involve crowdsourcing, participant recruitment, or direct human-subject experiments; it analyzes synthetic data and existing benchmarks.

\Needspace{8\baselineskip}

\item {\bf Institutional review board (IRB) approvals or equivalent for research with human subjects}
    \item[] Question: Does the paper describe potential risks incurred by study participants, whether such risks were disclosed to the subjects, and whether Institutional Review Board (IRB) approvals (or an equivalent approval/review based on the requirements of your country or institution) were obtained?
    \item[] Answer: \answerNA{}
    \item[] Justification: The research does not recruit or interact with human participants and conducts no human-subject study requiring participant-risk disclosure or IRB approval.

\Needspace{8\baselineskip}

\item {\bf Declaration of LLM usage}
    \item[] Question: Does the paper describe the usage of LLMs if it is an important, original, or non-standard component of the core methods in this research? Note that if the LLM is used only for writing, editing, or formatting purposes and does \emph{not} impact the core methodology, scientific rigor, or originality of the research, declaration is not required.
    %this research?
    \item[] Answer: \answerNA{}
    \item[] Justification: LLMs are not an important, original, or non-standard component of the core research methods. The fixed CLIP nuisance direction described in Appendix~\ref{app:hidden-slice-discovery} is a pretrained vision--language representation, not an LLM-based experimental method.

\end{enumerate}
